% answers_short.tex —— 2011 年·北京卷·理科 · 简短答案速查
% 由 tools/make_short_answers.py 自 answers.tex 的 \gksolhead/\gkanswer 提取，
% 并与 2011/2011_数学真题_北京理_答案要点.md 逐题比对后生成，勿手改（改源 MD 后重跑）。

\begin{gknotice}
\textbf{使用说明}：本节只给最终答案便于快速核对，完整推导见解析版。答案为编写者独立推导，\textbf{非官方标准答案}。
\end{gknotice}

\gkpart{选择题}
\begin{center}\begin{tabular}{ccccccccc}
\toprule
题号 & 1 & 2 & 3 & 4 & 5 & 6 & 7 & 8 \\
\midrule
答案 & C & A & B & D & A & D & C & C \\
\bottomrule
\end{tabular}\end{center}

\gkpart{填空题}
\begin{itemize}[leftmargin=2.6em]
  \item[\textbf{9.}] $\sin A=\dfrac{2\sqrt5}{5}$；$a=2\sqrt{10}$。
  \item[\textbf{10.}] $1$。
  \item[\textbf{11.}] $q=-2$；$|a_1|+|a_2|+\cdots+|a_n|=2^{n-1}-\dfrac12$。
  \item[\textbf{12.}] $14$。
  \item[\textbf{13.}] $(0,1)$。
  \item[\textbf{14.}] ②③。
\end{itemize}

\gkpart{解答题}
\begin{itemize}[leftmargin=2.6em]
  \item[\textbf{15.}] (1) 最小正周期 $T=\pi$；(2) 最大值 $2$（$x=\dfrac{\pi}{6}$），最小值 $-1$（$x=-\dfrac{\pi}{6}$）
  \item[\textbf{16.}] (1) 证明见解析（$BD\perp AC$，$BD\perp PA$）；(2) $\dfrac{\sqrt6}{4}$；(3) $PA=\sqrt6$
  \item[\textbf{17.}] (1) 平均数 $\dfrac{35}{4}$，方差 $\dfrac{11}{16}$；(2) $P(Y=17)=\dfrac18$，$P(Y=18)=P(Y=19)=P(Y=20)=\dfrac14$，$P(Y=21)=\dfrac18$；$E(Y)=19$
  \item[\textbf{18.}] (1) $k>0$ 时在 $(-\infty,-k)$、$(k,+\infty)$ 上单调递增，在 $(-k,k)$ 上单调递减；$k<0$ 时在 $(-\infty,k)$、$(-k,+\infty)$ 上单调递减，在 $(k,-k)$ 上单调递增；(2) $k\in\left[-\dfrac12,0\right)$
  \item[\textbf{19.}] (1) 焦点 $F_1(-\sqrt3,0)$、$F_2(\sqrt3,0)$，离心率 $e=\dfrac{\sqrt3}{2}$；(2) $|AB|=\dfrac{4\sqrt3\,|m|}{m^2+3}$（$|m|\geqslant1$），最大值 $2$（$m=\pm\sqrt3$ 时取得）
  \item[\textbf{20.}] (1) $A_5:0,1,0,1,0$（此时 $S(A_5)=2>0$）；(2) 证明见解析；递增时每步增 $1$，故 $a_n=12+1999=2011$，反之总和变化 $1999$ 只能由 $1999$ 个 $+1$ 实现；(3) 当且仅当 $n\equiv0,1\pmod 4$ 时存在；$n\equiv2,3\pmod 4$ 时不存在
\end{itemize}
